\subsection{A nonsplit additive division-ring extension: Problem 18.76}
\label{cand:18.76}

Let \(F=F(a,b)\), \(L=F(c,d)\), and \(G=F\times L\).
Free groups and their direct products admit two-sided
invariant orders. The Mal'cev--Neumann theorem therefore
embeds \(kG\) in a division ring for \(k=\Q\) or any
prime field; see \cite{sakarovitch,bergman}.
Let \(D\) be its division closure there. This is countable,
and \(G\leq D^\times\). Only elements from different
free factors are assumed to commute.

Define a derivation \(\delta:F\to kF\) by
\[
 \delta(xy)=\delta(x)+x\delta(y),\qquad
 \delta(a)=0,\quad\delta(b)=1.
\]
The prefix-sum formula on reduced words gives it
explicitly; contributions from adjacent inverse letters
cancel, with \(\delta(b^{-1})=-b^{-1}\).
In \(N=(D,+)\rtimes F\), with
\((r,x)(s,y)=(r+xs,xy)\), the maps
\[
 \gamma_c(r,x)=(cr,x),\qquad
 \gamma_d(r,x)=(dr+\delta(x),x)
\]
are automorphisms. The derivation law proves the
homomorphism identity; multiplication by \(c^{-1}\) and
the map \((r,x)\mapsto(d^{-1}(r-\delta(x)),x)\) give
their inverses. Freeness of \(L\) extends them to an
action, giving the actual group \(E=N\rtimes_\gamma L\).

Its quotient by \((D,+)\) is \(G=F\times L\), and the
conjugation action on the kernel is left multiplication
in \(D\). To check this and nonsplitting, define similarly
\(\eta:L\to kL\) by \(\eta(c)=0,\eta(d)=1\). Then
\[
 (r,x,h)(s,y,l)
   =(r+xhs+x\eta(h)\delta(y),\,xy,\,hl).
\]
This follows by composing the generator automorphisms;
the two derivations take values in commuting factor
algebras.

If a section existed, its values on \(a,b,c,d\) would
have arbitrary kernel coordinates \(u,v,w,z\in D\).
Put \(A_0=a-1,B_0=b-1,C_0=c-1,D_0=d-1\).
The four cross-factor commutations would imply
\[
 A_0w-C_0u=0,\quad A_0z-D_0u=0,\quad
 B_0w-C_0v=0,\quad B_0z-D_0v=1_D.
\]
Here \(1_D\) is a nonzero additive kernel element.
Since \(A_0,C_0\) are nonzero and invertible,
the first three equations give
\[
 w=A_0^{-1}C_0u,\qquad
 z=A_0^{-1}D_0u,\qquad
 v=B_0A_0^{-1}u.
\]
The left side of the fourth is then
\[
 B_0A_0^{-1}D_0u-D_0B_0A_0^{-1}u=0,
\]
because the two free factors commute, contradicting
\(1_D\ne0\). Thus the extension does not split, in any
of the stated characteristics. Its description as a
nested semidirect product does not split this particular
extension by \((D,+)\). Source:
\artifact{research/18.76-proof.md}.

